% Atur variabel berikut sesuai namanya

% nama
\newcommand{\name}{Ariq Maulana Tazakka}
\newcommand{\authorname}{Ariq, Maulana Tazakka}
\newcommand{\nickname}{Ariq}
\newcommand{\advisor}{Prof. Dr. I Ketut Eddy Purnama, S.T., M.T.}
\newcommand{\coadvisor}{Dr. Supeno Mardi Susiki Nugroho, S.T., M.T.}
\newcommand{\examinerone}{-Dr. Galileo Galilei, S.T., M.Sc}
\newcommand{\examinertwo}{-Friedrich Nietzsche, S.T., M.Sc}
\newcommand{\examinerthree}{-Alan Turing, ST., MT}
\newcommand{\headofdepartment}{Dr. Supeno Mardi Susiki Nugroho, S.T., M.T.}

% identitas
\newcommand{\nrp}{5024 21 1039}
\newcommand{\advisornip}{19690730 199512 1 001}
\newcommand{\coadvisornip}{19700313 199512 1 001}
\newcommand{\examineronenip}{-18560710 194301 1 001-}
\newcommand{\examinertwonip}{-18560710 194301 1 001-}
\newcommand{\examinerthreenip}{-18560710 194301 1 001-}
\newcommand{\headofdepartmentnip}{19700313199512 1 001} 

% judul
\newcommand{\tatitle}{VISUALISASI TEKANAN PADA PEMBULUH DARAH MENGGUNAKAN \emph{VIRTUAL REALITY} DAN ALGORITMA \emph{NETWORK FLOW} UNTUK SIMULASI SUMBATAN}
\newcommand{\engtatitle}{\emph{VISUALIZATION OF BLOOD VESSEL PRESSURE USING VIRTUAL REALITY AND NETWORK FLOW ALGORITHM FOR SIMULATING BLOCKAGES}}

% tempat
\newcommand{\place}{Surabaya}

% jurusan
\newcommand{\studyprogram}{Teknik Komputer}
\newcommand{\engstudyprogram}{Computer Engineering}

% fakultas
\newcommand{\faculty}{Teknologi Elektro dan Informatika Cerdas}
\newcommand{\engfaculty}{Intelligent Electrical and Informatics Technology}

% singkatan fakultas
\newcommand{\facultyshort}{FTEIC}
\newcommand{\engfacultyshort}{ELECTICS}

% departemen
\newcommand{\department}{Teknik Komputer}
\newcommand{\engdepartment}{Computer Engineering}

% kode mata kuliah
\newcommand{\coursecode}{EC234801}
